% Homework review for selected exercises from Chapters 1 and 2.
% Problem data follow homeworks.pdf. The supplied answers.pdf uses different
% numerical data for Section 2.1, Exercises 9 and 12; those two are solved
% directly from the textbook statements here.
\section[Homework]{Homework Solutions}
\sectioncard{Homework}{Selected Exercises from Chapters 1 and 2}

\begin{frame}{Homework Review: What We Will Practice}
\begin{columns}[T,onlytextwidth]
  \begin{column}{.48\textwidth}
    \textbf{Chapter 1}
    \begin{itemize}
      \item Linear independence
      \item Linear transformations
      \item Standard matrices
    \end{itemize}
  \end{column}
  \begin{column}{.48\textwidth}
    \textbf{Chapter 2}
    \begin{itemize}
      \item Matrix products
      \item Transposes and outer products
      \item Inverses and matrix equations
    \end{itemize}
  \end{column}
\end{columns}
\medskip
\begin{block}{A useful habit}
At every step, check dimensions and name the theorem or matrix operation being used.
\end{block}
\end{frame}

\begin{frame}{1.7, Exercise 2: Are the Vectors Independent?}
Determine whether
\[
\vect v_1=\begin{bmatrix}0\\0\\2\end{bmatrix},
\qquad
\vect v_2=\begin{bmatrix}0\\5\\-8\end{bmatrix},
\qquad
\vect v_3=\begin{bmatrix}-3\\4\\1\end{bmatrix}
\]
are linearly independent. Put the vectors into the columns of $A$:
\[
\left[ A\mid\vect0\right]
=\left[\begin{array}{ccc|c}
0&0&-3&0\\0&5&4&0\\2&-8&1&0
\end{array}\right]
\sim
\left[\begin{array}{ccc|c}
2&-8&1&0\\0&5&4&0\\0&0&-3&0
\end{array}\right].
\]
\pause
There is a pivot in every variable column, so $A\vect x=\vect0$ has only
$\vect x=\vect0$. Therefore
\[
\boxed{\{\vect v_1,\vect v_2,\vect v_3\}\text{ is linearly independent}.}
\]
\end{frame}

\begin{frame}{1.7, Exercise 8: Are the Columns Independent?}
\footnotesize
Let
\[
A=\begin{bmatrix}
1&-3&3&-2\\
-3&7&-1&2\\
0&1&-4&3
\end{bmatrix}.
\]
We solve $A\vect x=\vect0$:
\[
\begin{bmatrix}
1&-3&3&-2&0\\-3&7&-1&2&0\\0&1&-4&3&0
\end{bmatrix}
\xrightarrow{R_2\leftarrow R_2+3R_1}
\begin{bmatrix}
1&-3&3&-2&0\\0&-2&8&-4&0\\0&1&-4&3&0
\end{bmatrix}
\xrightarrow{R_2\leftarrow R_2+2R_3}
\begin{bmatrix}
1&-3&3&-2&0\\0&-2&8&-4&0\\0&0&0&2&0
\end{bmatrix}.
\]
Because there are at most three pivot positions yet there are four variables, the equation $Ax =
0$ has a \textbf{nontrivial solution} and the columns of A are linearly dependent.
\end{frame}

\begin{frame}{1.8, Exercise 3: Find a Preimage}
\footnotesize
For $T(\vect x)=A\vect x$, solve $A\vect x=\vect b$, where
\[
A=\begin{bmatrix}1&0&-2\\-2&1&6\\3&-2&-5\end{bmatrix},
\qquad
\vect b=\begin{bmatrix}-1\\7\\-3\end{bmatrix}.
\]
Row reduce the augmented matrix:
\[
\left[\begin{array}{ccc|c}
1&0&-2&-1\\-2&1&6&7\\3&-2&-5&-3
\end{array}\right]
\xrightarrow[R_3\leftarrow R_3-3R_1]{R_2\leftarrow R_2+2R_1}
\left[\begin{array}{ccc|c}
1&0&-2&-1\\0&1&2&5\\0&-2&1&0
\end{array}\right]
\xrightarrow{R_3\leftarrow R_3+2R_2}
\left[\begin{array}{ccc|c}
1&0&-2&-1\\0&1&2&5\\0&0&5&10
\end{array}\right].
\]
\pause
Back-substitution gives
$x_3=2$, $x_2=1$, and $x_1=3$.
Hence
\[
\boxed{\vect x=(3,1,2)^T}.
\]
Every variable is a pivot variable, so this preimage is unique.
\end{frame}

\begin{frame}{1.8, Exercise 9: Find the Kernel}
\footnotesize
Find all $\vect x\in\R^4$ satisfying $A\vect x=\vect0$ for
\[
A=\begin{bmatrix}1&-4&7&-5\\0&1&-4&3\\2&-6&6&-4\end{bmatrix}.
\]
\[
A\xrightarrow{R_3\leftarrow R_3-2R_1}
\begin{bmatrix}1&-4&7&-5\\0&1&-4&3\\0&2&-8&6\end{bmatrix}
\xrightarrow{R_3\leftarrow R_3-2R_2}
\begin{bmatrix}1&-4&7&-5\\0&1&-4&3\\0&0&0&0\end{bmatrix}.
\]
\pause
Let $x_3=s$ and $x_4=t$. Then
\[
x_2=4s-3t,
\qquad
x_1=9s-7t.
\]
Therefore
\[
\boxed{\vect x
=s\begin{bmatrix}9\\4\\1\\0\end{bmatrix}
+t\begin{bmatrix}-7\\-3\\0\\1\end{bmatrix}},
\qquad s,t\in\R.
\]
\end{frame}

\begin{frame}{1.8, Exercise 11: Is $\vect b$ in the Range?}
\footnotesize
Use the matrix from Exercise 9 and
\[
\vect b=\begin{bmatrix}-1\\1\\0\end{bmatrix}.
\]
The augmented matrix follows the same row operations:
\[
\left[\begin{array}{cccc|c}
1&-4&7&-5&-1\\0&1&-4&3&1\\2&-6&6&-4&0
\end{array}\right]
\longrightarrow
\left[\begin{array}{cccc|c}
1&-4&7&-5&-1\\0&1&-4&3&1\\0&2&-8&6&2
\end{array}\right]
\longrightarrow
\left[\begin{array}{cccc|c}
1&-4&7&-5&-1\\0&1&-4&3&1\\0&0&0&0&0
\end{array}\right].
\]
The system is consistent, so $\vect b$ is in the range if the transform.
\end{frame}

\begin{frame}{1.9, Exercise 6: Matrix of a Horizontal Shear}
The transformation leaves $\vect e_1$ unchanged and maps
$\vect e_2$ to $\vect e_2+5\vect e_1$.
\pause

The columns of the standard matrix are the images of the standard basis vectors:
\[
T(\vect e_1)=\begin{bmatrix}1\\0\end{bmatrix},
\qquad
T(\vect e_2)=\vect e_2+5\vect e_1
=\begin{bmatrix}5\\1\end{bmatrix}.
\]
Therefore
\[
\boxed{A=\begin{bmatrix}T(\vect e_1)&T(\vect e_2)\end{bmatrix}
=\begin{bmatrix}1&5\\0&1\end{bmatrix}}.
\]
Indeed, $A(x_1,x_2)^T=(x_1+5x_2,x_2)^T$, a horizontal shear.
\end{frame}

\begin{frame}{1.9, Exercise 15: Fill in the Matrix}
Find the matrix $A$ such that, for every $x_1,x_2,x_3$,
\[
A\begin{bmatrix}x_1\\x_2\\x_3\end{bmatrix}
=\begin{bmatrix}2x_1-3x_3\\4x_1\\x_1-x_2+x_3\end{bmatrix}.
\]
\pause
Read the coefficients in each output component as one row:
\[
\begin{array}{ccl}
2x_1+0x_2-3x_3 &\longrightarrow& \begin{bmatrix}2&0&-3\end{bmatrix},\\
4x_1+0x_2+0x_3 &\longrightarrow& \begin{bmatrix}4&0&0\end{bmatrix},\\
x_1-x_2+x_3 &\longrightarrow& \begin{bmatrix}1&-1&1\end{bmatrix}.
\end{array}
\]
Hence
\[
\boxed{A=\begin{bmatrix}2&0&-3\\4&0&0\\1&-1&1\end{bmatrix}}.
\]
\end{frame}

\begin{frame}{2.1, Exercise 9: When Does $AB=BA$?}
Let
\[
A=\begin{bmatrix}3&4\\-2&1\end{bmatrix},
\qquad
B=\begin{bmatrix}5&-6\\3&k\end{bmatrix}.
\]
Compute in both orders:
\[
AB=\begin{bmatrix}27&-18+4k\\-7&12+k\end{bmatrix},
\qquad
BA=\begin{bmatrix}27&14\\9-2k&12+k\end{bmatrix}.
\]
\pause
Equality requires corresponding entries to agree:
\[
-18+4k=14
\quad\text{and}\quad
-7=9-2k.
\]
Both equations give $\boxed{k=8}$. With this value, every entry of $AB$ and $BA$ agrees.
\end{frame}

\begin{frame}{2.1, Exercise 12: Build a Nonzero $B$ with $AB=0$}
Let
\[
A=\begin{bmatrix}2&-8\\-1&4\end{bmatrix}.
\]
Each column $\vect b$ of $B$ must satisfy $A\vect b=\vect0$.
Solving
\[
2x_1-8x_2=0
\quad\Longleftrightarrow\quad
x_1=4x_2
\]

Choose two different nonzero multiples:
\[
B=\begin{bmatrix}4&8\\1&2\end{bmatrix}.
\]
Then
\[
AB=
\begin{bmatrix}2&-8\\-1&4\end{bmatrix}
\begin{bmatrix}4&8\\1&2\end{bmatrix}
=\begin{bmatrix}0&0\\0&0\end{bmatrix}.
\]
\end{frame}

\begin{frame}{2.1, Exercise 35: Scalar Products}
Let
\[
\vect u=\begin{bmatrix}-2\\3\\-4\end{bmatrix},
\qquad
\vect v=\begin{bmatrix}a\\b\\c\end{bmatrix}.
\]
First compute the two $1\times1$ products:
\[
\vect u^T\vect v
=\begin{bmatrix}-2&3&-4\end{bmatrix}
\begin{bmatrix}a\\b\\c\end{bmatrix}
=-2a+3b-4c,
\]
\[
\vect v^T\vect u
=\begin{bmatrix}a&b&c\end{bmatrix}
\begin{bmatrix}-2\\3\\-4\end{bmatrix}
=-2a+3b-4c.
\]
Thus $\boxed{\vect u^T\vect v=\vect v^T\vect u}$.
\end{frame}

\begin{frame}{2.1, Exercise 35: Outer Products}
Now compute the two $3\times3$ products column by row:
\[
\vect u\vect v^T
=\begin{bmatrix}-2\\3\\-4\end{bmatrix}
\begin{bmatrix}a&b&c\end{bmatrix}
=\boxed{\begin{bmatrix}
-2a&-2b&-2c\\3a&3b&3c\\-4a&-4b&-4c
\end{bmatrix}},
\]
\pause
\[
\vect v\vect u^T
=\begin{bmatrix}a\\b\\c\end{bmatrix}
\begin{bmatrix}-2&3&-4\end{bmatrix}
=\boxed{\begin{bmatrix}
-2a&3a&-4a\\-2b&3b&-4b\\-2c&3c&-4c
\end{bmatrix}}.
\]
Notice that $\vect v\vect u^T=(\vect u\vect v^T)^T$; in general the two outer products are not equal.
\end{frame}

\begin{frame}{2.2, Exercise 3: Find the Inverse}
For
\[
A=\begin{bmatrix}8&3\\-7&-3\end{bmatrix},
\]
first compute the determinant:
\[
\det A=8(-3)-3(-7)=-24+21=-3\ne0.
\]
Therefore $A$ is invertible, and the $2\times2$ inverse formula gives
\[
A^{-1}
=-\frac{1}{3}\begin{bmatrix}-3&-3\\7&8\end{bmatrix}
=\boxed{\begin{bmatrix}1&1\\-\frac73&-\frac83\end{bmatrix}}.
\]
\pause
Quick check:
\[
A\,A^{-1}
=-\frac13\begin{bmatrix}8&3\\-7&-3\end{bmatrix}
\begin{bmatrix}-3&-3\\7&8\end{bmatrix}
=-\frac13\begin{bmatrix}-3&0\\0&-3\end{bmatrix}=I_2.
\]
\end{frame}

\begin{frame}{2.2, Exercise 16: True or False?}
\begin{block}{Statement}
If $A=\begin{bmatrix}a&b\\c&d\end{bmatrix}$ and $ad=bc$, then $A$ is not invertible.
\end{block}
\pause
The determinant is
\[
\det A=ad-bc.
\]
Under the condition $ad=bc$,
\[
\det A=ad-bc=0.
\]
\pause
A $2\times2$ matrix is invertible exactly when its determinant is nonzero.
Therefore the statement is \boxed{\text{True}}.
\end{frame}

\begin{frame}{2.2, Exercise 24: Cancel an Invertible Factor}
Suppose
\[
(B-C)D=0,
\]
where $B$ and $C$ are $m\times n$ matrices and $D$ is invertible.
Because the product is defined, $D$ is $n\times n$.
\pause

Multiply both sides on the \textbf{right} by $D^{-1}$:
\[
\bigl((B-C)D\bigr)D^{-1}=0D^{-1}.
\]
By associativity,
\[
(B-C)(DD^{-1})=0
\quad\Longrightarrow\quad
(B-C)I_n=0.
\]
Thus $B-C=0$, so
\[
\boxed{B=C}.
\]
The side matters: $D$ was on the right, so $D^{-1}$ must also be applied on the right.
\end{frame}

\begin{frame}{2.2, Exercise 29: Solve a Matrix Equation}
Assume $A$, $B$, and $C$ are invertible $n\times n$ matrices and solve
\[
C^{-1}(A+X)B^{-1}=I_n.
\]
\pause
Multiply on the left by $C$:
\[
(A+X)B^{-1}=C.
\]
Multiply on the right by $B$:
\[
A+X=CB.
\]
Subtract $A$:
\[
\boxed{X=CB-A}.
\]
\pause
Check:
\[
C^{-1}\bigl(A+(CB-A)\bigr)B^{-1}
=C^{-1}(CB)B^{-1}=I_n.
\]
\end{frame}

\begin{frame}{Homework Review: Ideas to Keep}
\begin{enumerate}
  \item To test independence, solve the homogeneous equation and look for free variables.
  \item To test whether $\vect b$ is in a range, test consistency of $A\vect x=\vect b$.
  \item The columns of a standard matrix are $T(\vect e_1),\ldots,T(\vect e_n)$.
  \item Matrix multiplication is order-sensitive; preserve the side of every operation.
  \item An invertible matrix can be cancelled by multiplying by its inverse on the correct side.
\end{enumerate}
\end{frame}
